\label{chap:83-18}

El límite inverso \(\Omega_\infty\) conserva la historia completa. Para
obtener una coordenada real hay que añadir una familia compatible de
intervalos y demostrar tres hechos diferentes: cada historia determina un
único valor; todo valor del compacto declarado es alcanzado; y las medidas
finitas descienden de forma consistente. Ninguno de estos hechos se deduce
del mero crecimiento del número de estados. La construcción que sigue
demuestra conjuntamente los tres hechos para el sistema HMT.

La publicación arquimediana conserva asimismo una aproximación de la
identidad por esperanzas condicionales cilíndricas. La concentración de masa
unitaria sobre prefijos compatibles proporciona el paso exacto desde los
cilindros nonádicos hasta la medida de Dirac, sin identificar densidad,
medida del soporte y masa total.

\begin{HMTScientificOwnerSubordinated}
\input{ampliaciones_sucesoras_20260824/parte_i_ii/owners/03_cilindros_dirac}
\end{HMTScientificOwnerSubordinated}

\section{Sistema de intervalos asociado a los estados finitos}

Sea \((S_n,r_{n,m})\) el sistema inverso de la
\cref{chap:sistema-inverso-continuidad-genealogica}. Para cada
\(a\in S_n\), asignamos un intervalo compacto no vacío
\begin{equation}
 I_n(a)=[\ell_n(a),u_n(a)]\subset\mathbb R.
 \label{p12-83-c18-s1:eq:intervalo-asociado-estado}
\end{equation}
La asignación es \emph{compatible con el refinamiento} cuando
\begin{equation}
 r_{n+1,n}(b)=a
 \quad\Longrightarrow\quad
 I_{n+1}(b)\subseteq I_n(a).
 \label{p12-83-c18-s1:eq:anidamiento-intervalos-estados}
\end{equation}
Es \emph{uniformemente contractiva} cuando existe una sucesión
\(\delta_n\downarrow0\) tal que
\begin{equation}
 \operatorname{diam}I_n(a)\leq\delta_n
 \qquad(a\in S_n).
 \label{p12-83-c18-s1:eq:contraccion-uniforme-intervalos}
\end{equation}

Para una historia \(x=(x_n)_n\), las inclusiones
\eqref{p12-83-c18-s1:eq:anidamiento-intervalos-estados} producen una cadena
\begin{equation}
 I_0(x_0)\supseteq I_1(x_1)\supseteq I_2(x_2)\supseteq\cdots.
 \label{p12-83-c18-s1:eq:cadena-intervalos-historia}
\end{equation}

\begin{theorem}[Evaluación de una historia]
\label{p12-83-c18-s1:thm:evaluacion-unica-historia-u024}
Si la asignación es compatible y uniformemente contractiva, entonces para
cada \(x\in\Omega_\infty\) existe un único número real
\begin{equation}
 \operatorname{ev}_{\mathbb R}(x)
 \in\bigcap_{n\geq0}I_n(x_n).
 \label{p12-83-c18-s1:eq:evaluacion-real-historia-u024}
\end{equation}
La aplicación
\begin{equation}
 \operatorname{ev}_{\mathbb R}:\Omega_\infty\longrightarrow\mathbb R
 \label{p12-83-c18-s1:eq:aplicacion-evaluacion-real-u024}
\end{equation}
es continua.
\end{theorem}

\begin{proof}
Los intervalos de \eqref{p12-83-c18-s1:eq:cadena-intervalos-historia} son compactos,
no vacíos y anidados. El teorema de los intervalos encajados garantiza que
su intersección es no vacía. Si \(s,t\) pertenecen a ella, entonces
\[
 |s-t|\leq\operatorname{diam}I_n(x_n)\leq\delta_n
\]
para todo \(n\). Al hacer \(n\to\infty\), se obtiene \(s=t\).

Para la continuidad, sea \(\varepsilon>0\) y elijamos \(n\) con
\(\delta_n<\varepsilon\). Si \(y\in[x_n]_n\), entonces
\(I_n(y_n)=I_n(x_n)\). Las dos evaluaciones pertenecen a ese intervalo y,
por tanto,
\[
 |\operatorname{ev}_{\mathbb R}(x)
  -\operatorname{ev}_{\mathbb R}(y)|<\varepsilon.
\]
El cilindro \([x_n]_n\) es un entorno abierto de \(x\).
\end{proof}

La construcción no identifica \(x\) con
\(\operatorname{ev}_{\mathbb R}(x)\). Una misma coordenada puede tener
varias historias, y una historia contiene más datos que cualquier punto de
la recta.

\section{Cobertura coherente del compacto visible}

Sea \(K\subset\mathbb R\) un compacto no vacío. La construcción HMT de los
cilindros asocia al anidamiento y a la contracción las identidades de
cobertura
\begin{align}
 K&=\bigcup_{a\in S_0}I_0(a),
 \label{p12-83-c18-s1:eq:cobertura-inicial-compacto-u024}\\
 I_n(a)&=
 \bigcup_{\substack{b\in S_{n+1}\\r_{n+1,n}(b)=a}}
 I_{n+1}(b)
 \qquad(a\in S_n).
 \label{p12-83-c18-s1:eq:cobertura-extensiones-compatible-u024}
\end{align}
La segunda igualdad es la publicación geométrica de la ley de prolongación:
cada punto de \(I_n(a)\) permanece en un cilindro hijo compatible. Junto con
la sobreyectividad de \(r_{n+1,n}\), ya establecida por la supervivencia de
los prefijos, produce simultáneamente estados por encima de \(a\) y la
cobertura íntegra de \(I_n(a)\).

\begin{theorem}[Sobreyectividad de la evaluación]
\label{p12-83-c18-s1:thm:sobreyectividad-evaluacion-compacto}
Las identidades construidas
\eqref{p12-83-c18-s1:eq:cobertura-inicial-compacto-u024} y
\eqref{p12-83-c18-s1:eq:cobertura-extensiones-compatible-u024} implican
\begin{equation}
 \operatorname{ev}_{\mathbb R}(\Omega_\infty)=K.
 \label{p12-83-c18-s1:eq:imagen-evaluacion-compacto}
\end{equation}
\end{theorem}

\begin{proof}
La inclusión de izquierda a derecha sigue de
\(I_n(x_n)\subseteq I_0(x_0)\subseteq K\).

Fijemos \(t\in K\). Definamos
\begin{equation}
 T_n(t):=\{a\in S_n:t\in I_n(a)\}.
 \label{p12-83-c18-s1:eq:estados-que-cubren-punto}
\end{equation}
La cobertura inicial hace no vacío \(T_0(t)\). Si \(a\in T_n(t)\), la
identidad \eqref{p12-83-c18-s1:eq:cobertura-extensiones-compatible-u024} proporciona al
menos un \(b\in T_{n+1}(t)\) con \(r_{n+1,n}(b)=a\). El grafo de estas
extensiones es finitamente ramificado y posee niveles no vacíos a toda
profundidad. Existe, por tanto, una cadena compatible \(x=(x_n)_n\) con
\(t\in I_n(x_n)\) para todo \(n\). La unicidad de la intersección obliga a
\(\operatorname{ev}_{\mathbb R}(x)=t\).
\end{proof}

\section{Cociente genealógico arquimediano y homeomorfismo con \texorpdfstring{\(\mathbb R\)}{R}}

Definimos la relación
\begin{equation}
 x\sim_{\mathbb R}y
 \quad\Longleftrightarrow\quad
 \operatorname{ev}_{\mathbb R}(x)
 =\operatorname{ev}_{\mathbb R}(y).
 \label{p12-83-c18-s1:eq:relacion-equivalencia-arquimediana}
\end{equation}
Sea
\begin{equation}
 q_{\mathbb R}:\Omega_\infty
 \longrightarrow\Omega_\infty/\!\sim_{\mathbb R}
 \label{p12-83-c18-s1:eq:proyeccion-cociente-arquimediano}
\end{equation}
la proyección canónica.

\begin{theorem}[Reconstrucción del compacto como cociente]
\label{p12-83-c18-s1:thm:homeomorfismo-cociente-compacto-u024}
Para el sistema HMT anterior existe un único homeomorfismo
\begin{equation}
 \overline{\operatorname{ev}}_{\mathbb R}:
 \Omega_\infty/\!\sim_{\mathbb R}
 \xrightarrow{\ \cong\ }K
 \label{p12-83-c18-s1:eq:homeomorfismo-cociente-compacto-u024}
\end{equation}
tal que
\begin{equation}
 \operatorname{ev}_{\mathbb R}
 =\overline{\operatorname{ev}}_{\mathbb R}\circ q_{\mathbb R}.
 \label{p12-83-c18-s1:eq:factorizacion-evaluacion-cociente}
\end{equation}
La fibra sobre \(t\in K\) es
\begin{equation}
 \mathfrak G_t
 :=\operatorname{ev}_{\mathbb R}^{-1}(t),
 \label{p12-83-c18-s1:eq:fibra-genealogias-mismo-valor}
\end{equation}
el conjunto completo de genealogías que publican \(t\).
\end{theorem}

\begin{proof}
La aplicación \(\operatorname{ev}_{\mathbb R}\) es continua, sobreyectiva y
constante exactamente en las clases de \(\sim_{\mathbb R}\). Por la
propiedad universal del cociente induce una aplicación continua y biyectiva
\(\overline{\operatorname{ev}}_{\mathbb R}\). El dominio es compacto por
ser imagen continua del compacto \(\Omega_\infty\), y \(K\) es Hausdorff.
Una biyección continua de un compacto a un Hausdorff es un homeomorfismo.
\end{proof}

El teorema prueba dos cosas y las mantiene separadas:
\begin{equation}
 \boxed{
 \begin{array}{c}
 \text{el espacio de historias es profinito y conserva memoria,}\\
 \text{su cociente por igualdad de evaluación es el compacto visible.}
 \end{array}}
 \label{p12-83-c18-s1:eq:dos-niveles-continuo-u024}
\end{equation}
No se afirma que \(\Omega_\infty\) sea homeomorfo a \(K\).

\section{Atlas genealógico de la recta real}

Para \(m\geq1\), sea \(K_m=[-m,m]\). La misma construcción, aplicada a cada
\(K_m\), produce un sistema \(\Omega_\infty^{(m)}\) que reconstruye
\(K_m\), junto con inclusiones cerradas
\begin{equation}
 j_m:\Omega_\infty^{(m)}\hookrightarrow
 \Omega_\infty^{(m+1)}
 \label{p12-83-c18-s1:eq:inclusion-atlas-genealogico}
\end{equation}
que satisfacen
\begin{equation}
 \operatorname{ev}_{m+1}\circ j_m
 =\iota_m\circ\operatorname{ev}_m,
 \label{p12-83-c18-s1:eq:compatibilidad-atlas-real}
\end{equation}
donde \(\iota_m:K_m\hookrightarrow K_{m+1}\) es la inclusión ordinaria.

\begin{proposition}[Agotamiento de la recta]
\label{p12-83-c18-s1:prop:agotamiento-recta-real-u024}
El cociente arquimediano del colímite estricto
\begin{equation}
 \Omega_\infty^{\mathbb R}
 :=\varinjlim_m\Omega_\infty^{(m)}
 \label{p12-83-c18-s1:eq:colimite-historias-recta}
\end{equation}
es
\begin{equation}
 \varinjlim_m[-m,m]
 =\bigcup_{m\geq1}[-m,m]
 =\mathbb R
 \label{p12-83-c18-s1:eq:colimite-compactos-recta}
\end{equation}
con su topología usual.
\end{proposition}

\begin{proof}
Cada pieza posee cociente \(K_m\) por el
\cref{p12-83-c18-s1:thm:homeomorfismo-cociente-compacto-u024};
\eqref{p12-83-c18-s1:eq:compatibilidad-atlas-real} entrelaza los cocientes. La topología
final del agotamiento por compactos coincide con la topología usual de
\(\mathbb R\): un subconjunto \(U\subseteq\mathbb R\) es abierto si y sólo
si \(U\cap[-m,m]\) es abierto en \([-m,m]\) para todo \(m\).
\end{proof}

El paso de un intervalo acotado a la recta no borra las fibras
\(\mathfrak G_t\). Las inclusiones \(j_m\) conservan la historia y sólo
amplían la carta visible disponible.

\section{Cartas ternaria balanceada y decimal por bloques}

La representación ternaria balanceada de un entero tiene la forma
\begin{equation}
 n=\sum_{j=0}^{J}\epsilon_j3^j,
 \qquad
 \epsilon_j\in\{-1,0,+1\},
 \label{p12-83-c18-s1:eq:carta-ternaria-balanceada-u024}
\end{equation}
con una regla de acarreo fijada. Esta carta conserva el signo local y el
cociente de la reducción; no debe confundirse con la palabra no negativa en
\(\mathbb F_3\).

La carta decimal por bloques utiliza
\begin{equation}
 D_N=1000D_{N-1}+u_N,
 \qquad
 u_N\in\{0,\ldots,999\},
 \qquad D_0=0.
 \label{p12-83-c18-s1:eq:concatenacion-base-mil-u024}
\end{equation}
El bloque \(u_N\), su clase módulo nueve y el acarreo son campos distintos.
Como \(1000\equiv1\pmod9\),
\begin{equation}
 D_N\equiv\sum_{j=1}^{N}u_j\pmod9.
 \label{p12-83-c18-s1:eq:base-mil-compatible-mod9-u024}
\end{equation}
Esta congruencia no recupera el orden de los bloques; la secuencia
\((u_1,\ldots,u_N)\) permanece en el registro.

Para \(M\geq1\), las celdas decimales semiabiertas son
\begin{equation}
 \mathcal D_{M,q}
 :=[q10^{-M},(q+1)10^{-M}),
 \qquad q\in\mathbb Z.
 \label{p12-83-c18-s1:eq:celda-decimal-semiabierta-u024}
\end{equation}
Un intervalo semiabierto \([\ell,h)\) está contenido en una única celda si
y sólo si
\begin{equation}
 \left\lfloor10^M\ell\right\rfloor
 =\left\lceil10^Mh\right\rceil-1.
 \label{p12-83-c18-s1:eq:criterio-exacto-celda-decimal-u024}
\end{equation}
La fórmula con techo en el extremo derecho cubre también el caso en que
\(h\) coincide exactamente con una frontera decimal.

\begin{definition}[Truncamiento y redondeo]
Para \(x\geq0\), definimos
\begin{align}
 \operatorname{Tr}_M(x)
 &:=10^{-M}\left\lfloor10^Mx\right\rfloor,
 \label{p12-83-c18-s1:eq:truncamiento-decimal-u024}\\
 \operatorname{Rd}_M(x)
 &:=10^{-M}\left\lfloor10^Mx+\frac12\right\rfloor.
 \label{p12-83-c18-s1:eq:redondeo-decimal-u024}
\end{align}
\end{definition}
Estas operaciones coinciden salvo cuando la fracción descartada cruza el
umbral de media unidad en la última posición. La distinción es semántica y
aritmética; no puede sustituirse una por otra para forzar compatibilidad de
prefijos.

\section{Descenso de operaciones por las fibras de evaluación}

Sea \(D_\star\subseteq\Omega_\infty\times\Omega_\infty\) el dominio de una
operación parcial continua
\begin{equation}
 \star:D_\star\longrightarrow\Omega_\infty.
 \label{p12-83-c18-s1:eq:operacion-historias-u024}
\end{equation}
Suponemos que \(D_\star\) es saturado por las fibras de
\(\operatorname{ev}_{\mathbb R}\times\operatorname{ev}_{\mathbb R}\).

\begin{theorem}[Criterio exacto de descenso]
\label{p12-83-c18-s1:thm:criterio-descenso-operaciones-u024}
Existe una única operación visible
\begin{equation}
 \overline\star:
 (\operatorname{ev}_{\mathbb R}\times\operatorname{ev}_{\mathbb R})
 (D_\star)\longrightarrow K
 \label{p12-83-c18-s1:eq:operacion-descendida-visible}
\end{equation}
que satisface
\begin{equation}
 \operatorname{ev}_{\mathbb R}(x\star y)
 =\operatorname{ev}_{\mathbb R}(x)
  \ \overline\star\ 
  \operatorname{ev}_{\mathbb R}(y)
 \label{p12-83-c18-s1:eq:entrelazamiento-operacion-descendida}
\end{equation}
si y sólo si
\begin{equation}
 \begin{split}
 &\operatorname{ev}_{\mathbb R}(x)=
  \operatorname{ev}_{\mathbb R}(x'),\\
 &\operatorname{ev}_{\mathbb R}(y)=
  \operatorname{ev}_{\mathbb R}(y')
 \end{split}
 \quad\Longrightarrow\quad
 \operatorname{ev}_{\mathbb R}(x\star y)=
 \operatorname{ev}_{\mathbb R}(x'\star y')
 \label{p12-83-c18-s1:eq:constancia-operacion-sobre-fibras}
\end{equation}
para todos los pares del dominio.
\end{theorem}

\begin{proof}
Si \(\overline\star\) existe, el miembro derecho de
\eqref{p12-83-c18-s1:eq:entrelazamiento-operacion-descendida} depende sólo de las dos
evaluaciones, y se obtiene
\eqref{p12-83-c18-s1:eq:constancia-operacion-sobre-fibras}. Recíprocamente, esta
constancia permite definir
\[
 s\ \overline\star\ t
 :=\operatorname{ev}_{\mathbb R}(x\star y)
\]
para cualesquiera representantes \(x,y\) con evaluaciones \(s,t\). La
condición prueba que la definición no depende de los representantes. La
unicidad es inmediata.
\end{proof}

Las hojas aditiva y multiplicativa de APP suministran operaciones en un
soporte común, pero sus acarreos y sus memorias no se identifican. El
teorema anterior permite que suma y producto desciendan a la carta visible
sin exigir que sus realizaciones genealógicas coincidan.

\section{Correspondencia entre cardinalidad discreta, medida cilíndrica y publicación arquimediana}

En cada profundidad aparecen tres operaciones que deben conservarse
separadas:
\begin{align}
 \operatorname{Count}_n(A)&:=|A|,
 \label{p12-83-c18-s1:eq:operacion-conteo-u024}\\
 \operatorname{Measure}_n(A)&:=\mu_n(A),
 \label{p12-83-c18-s1:eq:operacion-medida-u024}\\
 \operatorname{Publish}_n(a)&:=L_n(a).
 \label{p12-83-c18-s1:eq:operacion-publicacion-u024}
\end{align}
El conteo depende del número de estados; la medida depende de pesos
compatibles; la publicación depende del lector. Dos fibras de igual
cardinal pueden tener pesos distintos, y dos estados de igual publicación
pueden pertenecer a cilindros distintos. La igualdad de una de estas tres
salidas no implica la igualdad de las otras dos.

\section{Obstrucciones de la evaluación arquimediana: anidamiento, continuidad, cociente, medidas proyectivas y descenso de operaciones}

\begin{criterion}[Refutación de la evaluación y de la medida]
La construcción queda refutada si se verifica cualquiera de los hechos
siguientes:
\begin{enumerate}
 \item una extensión compatible recibe un intervalo que no está contenido
       en el intervalo de su restricción;
 \item los diámetros no tienden a cero pero se afirma unicidad de la
       intersección;
 \item la evaluación no es continua en la topología de cilindros;
 \item se afirma que todo punto de \(K\) aparece sin demostrar la cobertura
       coherente \eqref{p12-83-c18-s1:eq:cobertura-extensiones-compatible-u024};
 \item la aplicación inducida por el cociente no es biyectiva;
 \item se identifica el espacio profinito con su cociente visible;
 \item las inclusiones del atlas no entrelazan sus evaluaciones;
 \item se usa \(\lfloor10^Mh\rfloor\) en lugar de
       \(\lceil10^Mh\rceil-1\) para un intervalo semiabierto y se pierde
       el caso de frontera;
 \item se confunden truncamiento y redondeo;
 \item una operación descendida depende del representante elegido en una
       fibra;
 \item una familia \((\mu_n)_n\) viola
       \eqref{eq:consistencia-proyectiva-medidas-u024};
 \item la medida de un cilindro depende del nivel al que se lo refina;
 \item una medida condicionada deja de sumar uno o pierde consistencia;
 \item se identifica la medida visible \(\nu\) con las medidas internas
       \(\mu_t\);
 \item se deduce una flecha entre \(\mathbb R\) y
       \(\mathcal E_\infty\) sólo porque ambas publicaciones comparten
       antecedente;
 \item se usa el conteo de una fibra como si determinara su medida o su
       valor publicado.
\end{enumerate}
\end{criterion}


\section{Del espacio de historias a la recta real}

Sea
\[
 \Psi:\mathbb Z_3^6\xrightarrow{\ \sim\ }
       (\mathbb F_3^6)^{\mathbb N}
\]
la conjugación Hensel entre la memoria inicial y la cinta completa de
bloques de seis trits. Concatenar las seis coordenadas de cada bloque da
una cinta ternaria
\(
 (a_n)_{n\geq 1}\in\{0,1,2\}^{\mathbb N}
\), sobre la que actúa la evaluación
\[
 \operatorname{ev}_3((a_n))
 :=\sum_{n\geq1}\frac{a_n}{3^n}\in[0,1].
\]
Las dos expansiones de un extremo ternario representan el mismo valor y
permanecen diferenciadas antes de aplicar la evaluación. Esto no constituye
una ambigüedad del objeto genealógico: es precisamente una fibra del lector.

Definimos
\[
 X_{\mathrm{raw}}:=\mathbb Z\times\mathbb Z_3^6,
 \qquad
 P(m,x):=m+\operatorname{ev}_3(\operatorname{flat}(\Psi x)).
\]
La coordenada entera proporciona el atlas numerable de la recta y la
coordenada 3-ádica conserva la historia completa de la parte fraccionaria.

\begin{theorem}[Reconstrucción exacta del cociente arquimediano]
\label{p12-83-c18-s2:hmtch:thm-cociente-real}
La aplicación \(P:X_{\mathrm{raw}}\to\mathbb R\) es sobreyectiva. Si
\[
 x\sim_P y
 \quad\Longleftrightarrow\quad
 P(x)=P(y),
\]
entonces
\[
 \overline P:X_{\mathrm{raw}}/\!\sim_P\;\longrightarrow\mathbb R,
 \qquad [x]\longmapsto P(x),
\]
es una biyección. En consecuencia,
\[
 \left|X_{\mathrm{raw}}/\!\sim_P\right|
 =|\mathbb R|
 =2^{\aleph_0}.
\]
\end{theorem}

\begin{proof}
Toda sucesión ternaria se agrupa de manera única en bloques consecutivos de
seis cifras. La conjugación \(\Psi\) asigna a esa cinta una única memoria
3-ádica inicial. Dado \(r\in\mathbb R\), se toma
\(m=\lfloor r\rfloor\) y una expansión ternaria de
\(r-m\in[0,1)\); la reconstrucción anterior produce
\(x\in\mathbb Z_3^6\) con \(P(m,x)=r\). Así, \(P\) es sobreyectiva.

La aplicación \(\overline P\) está bien definida por la definición de
\(\sim_P\), es sobreyectiva porque lo es \(P\), e inyectiva porque dos
clases con la misma imagen pertenecen a una misma fibra. Finalmente,
\[
 |\mathbb Z_3^6|
 =\left|(\mathbb F_3^6)^{\mathbb N}\right|
 =729^{\aleph_0}
 =2^{\aleph_0},
\]
y multiplicar por la coordenada numerable \(\mathbb Z\) no altera ese
cardinal. El cálculo coincide con el del cociente, ya identificado con
\(\mathbb R\).
\end{proof}

\begin{remark}[Clausura superviviente]
\label{p12-83-c18-s2:hmtch:rem-version-superviviente}
La completación bruta y el subespacio superviviente
\(X_{\mathrm{surv}}\) son dos etapas de una misma construcción. La ley de
supervivencia produce coberturas coherentes en cada nivel, contracción
uniforme de los diámetros y una prolongación para toda celda visible. El lema
de la rama infinita produce entonces una sección sobre cada punto, y la
evaluación induce, compacto a compacto, un homeomorfismo del cociente con el
intervalo publicado. El agotamiento por \([-m,m]\) reconstruye toda la recta,
en concordancia con los
\cref{p12-83-c18-s1:prop:agotamiento-recta-real-u024,p12-83-c18-s2:hmtch:thm-cociente-real}.
\end{remark}

\section{Media arquimediana y doble proyección del mismo estado}

Sea \(u=(m,\mathbf u)\in
\mathscr O_{\mathbb R,\leftrightarrow}^{\rm HMT}\). Un prefijo
\(\mathbf u_n\) determina en la carta fraccionaria un cilindro orientado
semiabierto
\[
 I_n(\mathbf u_n)=[a_n,b_n),
\]
con la convención complementaria en los extremos de doble expansión. Para
aplicar compacidad se usa su clausura
\(K_n=[a_n,b_n]\). Las compatibilidades dan
\[
 K_{n+1}\subseteq K_n,
 \qquad
 \delta_n:=b_n-a_n\longrightarrow0.
\]
La media arquimediana global del cilindro es
\[
 \mu_n(u)=m+\frac{a_n+b_n}{2}.
\]
Si \(q\geq n\), entonces
\[
 |\mu_q(u)-\mu_n(u)|\leq\delta_n,
\]
de modo que
\begin{equation}
 \mathfrak R(u):=\lim_{n\to\infty}\mu_n(u)
 \label{p12-83-c18-s2:hmtch:eq-media-arquimediana}
\end{equation}
existe. Lo que tiende a cero es el diámetro del cilindro; el valor real es
la resultante del encajamiento, no aquello que se anula.

\begin{theorem}[Cociente two-way y recta real]
\label{p12-83-c18-s2:hmtch:thm-cociente-two-way-real}
Sea
\[
 u\sim_{\mathfrak R}v
 \quad\Longleftrightarrow\quad
 \mathfrak R(u)=\mathfrak R(v).
\]
Bajo la cobertura coherente de las cartas HMT, la evaluación induce un
homeomorfismo
\begin{equation}
 \overline{\mathfrak R}:
 \mathscr O_{\mathbb R,\leftrightarrow}^{\rm HMT}/\!
 \sim_{\mathfrak R}
 \xrightarrow{\ \sim\ }\mathbb R.
 \label{p12-83-c18-s2:hmtch:eq-homeomorfismo-two-way-real}
\end{equation}
En particular, éste es un morfismo HMT continuo: no el resultado de elegir
una historia aislada de cada fibra.
\end{theorem}

\begin{proof}
La cota de cilindros prueba la continuidad de \(\mathfrak R\). La cobertura
produce una historia sobre cada real y, por tanto, sobreyectividad. Al pasar
al cociente por las fibras, la aplicación inducida es biyectiva. En cada
carta compacta, una biyección continua hacia un intervalo de Hausdorff es un
homeomorfismo. La topología global del cociente es la topología final del
pegado de esas cartas; la fibra de evaluación identifica el extremo
\(m+1\) de una carta con el extremo \(0\) trasladado de la carta siguiente.
El recubrimiento resultante es localmente finito, por lo que los
homeomorfismos de carta se pegan en un homeomorfismo global hacia
\(\mathbb R\).
\end{proof}
